% problem_id: S001-theorem-multiplicity-with-koszul
% source_id: S001 (Stacks Project)
% source_file: intersection.tex
% statement_locator: intersection.tex lines 1001--1012
% proof_locator: intersection.tex lines 1014--1089
% env: theorem
% id_note: label 在本来源内唯一
% status: complete (statement + author proof verbatim + reason + locators)
%
\begin{theorem}
\label{theorem-multiplicity-with-koszul}
\begin{reference}
\cite[Theorem 1 in part B of Chapter IV]{Serre_algebre_locale}
\end{reference}
Let $A$ be a Noetherian local ring. Let $I = (f_1, \ldots, f_r) \subset A$
be an ideal of definition. Let $M$ be a finite $A$-module. Then
$$
e_I(M, r) = \sum
(-1)^i\text{length}_A H_i(K_\bullet(f_1, \ldots, f_r) \otimes_A M)
$$
\end{theorem}

\begin{proof}
Let us change the Koszul complex $K_\bullet(f_1, \ldots, f_r)$ into a cochain
complex $K^\bullet$ by setting $K^n = K_{-n}(f_1, \ldots, f_r)$.
Then $K^\bullet$ is sitting in degrees $-r, \ldots, 0$ and
$H^i(K^\bullet \otimes_A M) = H_{-i}(K_\bullet(f_1, \ldots, f_r) \otimes_A M)$.
The statement of the theorem makes sense as the modules
$H^i(K^\bullet \otimes M)$ are annihilated by $f_1, \ldots, f_r$
(More on Algebra, Lemma \ref{more-algebra-lemma-homotopy-koszul})
hence have finite length.
Define a filtration on the complex $K^\bullet$ by setting
$$
F^p(K^n \otimes_A M) =
I^{\max(0, p + n)}(K^n \otimes_A M),\quad p \in \mathbf{Z}
$$
Since $f_i I^p \subset I^{p + 1}$ this is a filtration by subcomplexes.
Thus we have a filtered complex and we obtain a spectral sequence, see
Homology, Section \ref{homology-section-filtered-complex}.
We have
$$
E_0 = \bigoplus\nolimits_{p, q} E_0^{p, q} =
\bigoplus\nolimits_{p, q} \text{gr}^p(K^{p + q} \otimes_A M) =
\text{Gr}_I(K^\bullet \otimes_A M)
$$
Since $K^n$ is finite free we have
$$
\text{Gr}_I(K^\bullet \otimes_A M) =
\text{Gr}_I(K^\bullet) \otimes_{\text{Gr}_I(A)} \text{Gr}_I(M)
$$
Note that $\text{Gr}_I(K^\bullet)$ is the Koszul
complex over $\text{Gr}_I(A)$ on the elements
$\overline{f}_1, \ldots, \overline{f}_r \in I/I^2$.
A simple calculation (omitted)
shows that the differential $d_0$ on $E_0$
agrees with the differential coming from the Koszul complex.
Since $\text{Gr}_I(M)$ is a finite $\text{Gr}_I(A)$-module
and since $\text{Gr}_I(A)$ is Noetherian (as a quotient
of $A/I[x_1, \ldots, x_r]$ with $x_i \mapsto \overline{f}_i$), the
cohomology module $E_1 = \bigoplus E_1^{p, q}$
is a finite $\text{Gr}_I(A)$-module. However, as above
$E_1$ is annihilated by $\overline{f}_1, \ldots, \overline{f}_r$.
We conclude $E_1$ has finite length.
In particular we find that $\text{Gr}^p_F(K^\bullet \otimes M)$ is
acyclic for $p \gg 0$.

\medskip\noindent
Next, we check that the spectral sequence above converges
using Homology, Lemma \ref{homology-lemma-filtered-complex-ss-converges}.
The required equalities follow easily from the Artin-Rees lemma
in the form stated in Algebra, Lemma \ref{algebra-lemma-map-AR}.
Thus we see that
\begin{align*}
\sum (-1)^i\text{length}_A(H^i(K^\bullet \otimes_A M))
& =
\sum (-1)^{p + q} \text{length}_A(E_\infty^{p, q}) \\
& =
\sum (-1)^{p + q} \text{length}_A(E_1^{p, q})
\end{align*}
because as we've seen above the length of $E_1$ is finite
(of course this uses additivity of lengths). Pick $t$ so
large that $\text{Gr}^p_F(K^\bullet \otimes M)$
is acyclic for $p \geq t$ (see above). Using
additivity again we see that
$$
\sum (-1)^{p + q} \text{length}_A(E_1^{p, q}) =
\sum\nolimits_n \sum\nolimits_{p \leq t}
(-1)^n \text{length}_A(\text{gr}^p(K^n \otimes_A M))
$$
This is equal to
$$
\sum\nolimits_{n = -r, \ldots, 0} (-1)^n{r \choose |n|} \chi_{I, M}(t + n)
$$
by our choice of filtration above and the definition of $\chi_{I, M}$ in
Algebra, Section \ref{algebra-section-Noetherian-local}.
The lemma follows from Lemma \ref{lemma-leading-coefficient}
and the definition of $e_I(M, r)$.
\end{proof}
